% Abstract
Predicting cellular commitment thresholds---such as ATP concentrations triggering irreversible developmental transitions---requires modeling regulatory signal consumption during decision-making. Biological Petri net theory lacks formal semantics for such hierarchical regulatory signals. Classical test arcs model catalysis (non-consuming reads) but cannot express signal consumption essential for commitment finality and pathway preemption. We introduce \textit{signal hierarchy theory} providing signal places ($\Psi$) as information channels distinct from metabolic mass transfer, with signal flow arcs enabling consumption semantics. Hierarchical layer assignment formalizes multi-scale organization where higher regulatory layers control lower metabolic layers. We prove layer assignment consistency for acyclic signal networks and preemption correctness guarantees. The formalism bridges the expressiveness gap between metabolic mass transfer (normal arcs) and regulatory information flow (signal flow arcs), enabling unified modeling where energy status functions simultaneously as metabolite and regulatory signal. Application to \textit{Bacillus subtilis} sporulation demonstrates quantitative capability: signal hierarchy analysis computes ATP commitment threshold (2.38 mM) with 7\% accuracy versus experimental measurements, a prediction impossible in classical Bio-PN lacking consumption semantics. The validated framework enables quantitative threshold determination, basin geometry quantification, and commitment reliability analysis previously inexpressible in Bio-PN theory.
